% ----------------------------------------------------------------------
% Chapter 2: The setting -- discrete-time market with frictions and
% convex risk measures.
% Standalone compilable artifact (own preamble); when assembling the
% thesis, strip the preamble and merge the body only.
% Counterpart of Sections 2 and 3 of Buehler et al. (2019), but on a
% general probability space.
% Cross-references to other chapters are written as plain text
% ("Chapter 3", "Chapter 4"); these are stable (chapters keep their
% numbers), so no FAKE REF replacement is needed for them.
% ----------------------------------------------------------------------
\documentclass[11pt,a4paper]{report}
\usepackage[T1]{fontenc}
\usepackage[utf8]{inputenc}
\usepackage{amsmath,amssymb,amsthm}
\usepackage{enumitem}

\numberwithin{equation}{chapter}

\theoremstyle{plain}
\newtheorem{proposition}{Proposition}[chapter]
\newtheorem{lemma}[proposition]{Lemma}
\newtheorem{theorem}[proposition]{Theorem}
\newtheorem{corollary}[proposition]{Corollary}
\theoremstyle{definition}
\newtheorem{assumption}[proposition]{Assumption}
\newtheorem{definition}[proposition]{Definition}
\newtheorem{example}[proposition]{Example}
% Upright body, bold head (not the italic amsthm "remark" style).
\newtheorem{remark}[proposition]{Remark}

\newcommand{\E}{\mathbb{E}}
\newcommand{\PP}{\mathbb{P}}
\newcommand{\R}{\mathbb{R}}
\newcommand{\N}{\mathbb{N}}
\newcommand{\F}{\mathcal{F}}
\newcommand{\HH}{\mathcal{H}}
\newcommand{\Hu}{\mathcal{H}^{\mathrm{u}}}
\newcommand{\as}{\ \PP\text{-a.s.}}

\begin{document}

\setcounter{chapter}{1}

\chapter{The setting: a discrete-time market with frictions and convex risk measures}

This chapter sets up all objects needed in the rest of the thesis. We first construct the market: trading dates, information, hedging instruments, trading strategies, constraints and transaction costs. Then we use these objects to define the terminal profit-and-loss equation. The second half of the chapter develops convex risk measures, and in particular the class of optimized certainty equivalents, to measure the ``goodness'' of the P\&L.

The market model (Sections~\ref{sec:market} to~\ref{sec:pnl}) copies the setup of B\"uhler et al.\ \cite[Section~2]{buehler}, with one structural difference: the probability space is \emph{general}, not finite.

\section{The market}\label{sec:market}

Fix a probability space $(\Omega,\F,\PP)$. We emphasize that $(\Omega,\F,\PP)$ is a general probability space: in contrast to B\"uhler et al.\ \cite{buehler}, we do not assume that $\Omega$ is finite. All identities and inequalities between random variables are understood in the $\PP$-almost sure sense.

Trading takes place at finitely many dates
\[
    0=t_0<t_1<\dots<t_n=T.
\]
New market information available at $t_k$ is modelled by an $\R^r$-valued random vector $I_k$; the vector $I_k$ may contain market prices of the hedging instruments, but also additional signals such as prices of non-traded instruments, news or trading volumes. The flow of information is described by the filtration $\mathbb{F}=(\F_k)_{k=0,\dots,n}$ generated by the \emph{information process} $I=(I_k)_{k=0,\dots,n}$,
\[
    \F_k:=\sigma(I_0,\dots,I_k),\qquad k=0,\dots,n,
\]
and we set $\F_T:=\F_n$. Without loss of generality we take $\F=\F_T$.

The market contains $d$ hedging instruments whose \emph{mid-prices} are given by an $\R^d$-valued $\mathbb{F}$-adapted process $S=(S_k)_{k=0,\dots,n}$, where $S_k$ is the price vector at date $t_k$. The instruments may be primary assets such as stocks or currencies, but also secondary assets such as liquidly traded options; this flexibility is one of the strengths of the framework, because hedging a complex derivative portfolio in practice relies on liquid vanilla options as much as on the underlying itself. We stress that $S$ is simply an adapted process: no martingale-measure assumption of any kind is made at this point.

The \emph{liability} is an $\F_T$-measurable random variable $Z$, interpreted as the payoff of a portfolio of derivatives that the trader has sold and must service at time $T$. As in B\"uhler et al.\ \cite[Section~2]{buehler}, we adopt the following simplifications. Interest rates are zero, and any cash flow before $T$ is assumed to be invested into the overnight risk-free account. Under zero rates, this is costless, so we can assume all payments are due at the terminal date $T$. We also exclude American-style optionality.

\section{Trading strategies, costs and the P\&L equation}\label{sec:pnl}

A \emph{trading strategy} is an $\R^d$-valued $\mathbb{F}$-adapted process $\delta=(\delta_k)_{k=0,\dots,n-1}$, where $\delta_k^i$ is the position in the $i$-th hedging instrument held over the interval $(t_k,t_{k+1}]$. Adaptedness means that the position chosen at $t_k$ may only use information available at $t_k$. We use throughout the convention
\[
    \delta_{-1}:=0,\qquad\delta_n:=0:
\]
the trader starts with no position and liquidates the entire portfolio at $T$. We write $\Hu$ for the set of all such strategies (the superscript stands for ``unconstrained'').

In real markets, not every position is available: instruments may be tradeable only from certain dates onwards, liquidity limits the position size and risk management imposes bounds or short-selling restrictions. Following B\"uhler et al.\ \cite[Section~2.4]{buehler}, such \emph{trading constraints} are encoded by continuous maps $H_k$ that project an intended position onto the set of admissible ones, in a way that depends on the information available up to $t_k$ and on the previously held position $\delta_{k-1}$. The constrained version of $\delta\in\Hu$ is denoted by $H\circ\delta$ and defined recursively by
\[
    (H\circ\delta)_k:=H_k\bigl((H\circ\delta)_0,\dots,(H\circ\delta)_{k-1},\delta_k\bigr),\qquad k=0,\dots,n-1,
\]
with $H_k(\cdot,0)=0$, so that the zero strategy is always admissible. The set of admissible strategies is
\[
    \HH:=\{H\circ\delta:\delta\in\Hu\}.
\]
In Chapters~3 and~4 we will restrict admissibility by requiring that strategies cannot trade into a loss.

Trading is costly. For each $k$, the \emph{cost function} $c_k\colon\Omega\times\R^d\to[0,\infty)$ gives the cost $c_k(x)$ of changing the position by $x\in\R^d$ at date $t_k$; it is $\F_k\otimes\mathcal{B}(\R^d)$-measurable, upper semicontinuous in the spatial variable and normalized so that
\[
    c_k(0)=0:
\]
not trading is free and costs are nonnegative. The \emph{total cost} of a strategy $\delta$ is
\[
    C_T(\delta):=\sum_{k=0}^{n}c_k(\delta_k-\delta_{k-1}),
\]
where the conventions $\delta_{-1}=\delta_n=0$ ensure that the initial build-up and the final liquidation of the position are charged as well.


\begin{example}\label{ex:prop-costs}
    Following B\"uhler et al.\ \cite[Section~2]{buehler}, we give an example of proportional
    transaction costs that given by
    \[
        c_k(x)=\sum_{i=1}^d\kappa^i S_k^i|x^i|
    \]
    with relative cost parameters $\kappa^i\ge0$: each trade costs a fixed proportion of the transacted volume. These cost functions are continuous, convex and normalized.
\end{example}

\noindent Upper semicontinuity of the cost functions is exactly what the limit arguments of Chapter~4 require: along a convergent sequence of strategies, costs may drop in the limit but not jump up.

The trading gains of $\delta$ from price movements are given by the discrete stochastic integral
\[
    (\delta\cdot S)_T:=\sum_{k=0}^{n-1}\delta_k\cdot(S_{k+1}-S_k).
\]
Now suppose the trader sells the liability $Z$ for an initial cash amount $p_0$ and trades according to $\delta\in\HH$. Collecting all cash flows at $T$ yields the \emph{terminal profit-and-loss} of B\"uhler et al.\ \cite[Equation~(1)]{buehler},
\begin{equation}\label{eq:PL}
    \mathrm{PL}_T(Z,p_0,\delta):=-Z+p_0+(\delta\cdot S)_T-C_T(\delta).
\end{equation}
Equation \eqref{eq:PL} is the central object of this thesis. Its four terms are the payment of the liability, the premium received, the trading gains and the frictional losses. The remainder of the thesis is organized around the question of how to determine the riskiness of the random variable \eqref{eq:PL}, how to choose $\delta$ (Chapter~3) and to compute such a $\delta$ (Chapter~4) so as to minimize its risk as much as possible.

\section{Relevant markets and feasibility}\label{sec:arbitrage}

It is convenient to isolate the part of the P\&L that is generated by pure trading. For $\delta\in\HH$, define the \emph{gain} of $\delta$ net of costs by
\[
    W(\delta):=(\delta\cdot S)_T-C_T(\delta).
\]

A first natural question is whether the market allows arbitrage: a strategy $\delta\in\HH$ with
\[
    W(\delta)\ge0\as\qquad\text{and}\qquad\PP[W(\delta)>0]>0.
\]
Classical pricing theory begins by excluding such strategies. The deep hedging framework does not: no no-arbitrage assumption is imposed a priori. This is a deliberate modelling choice. The framework is designed to operate on market data or on scenarios from a simulator, and neither is guaranteed to be arbitrage-free. What must be avoided is not arbitrage itself but a market that is \emph{too} attractive, in the following sense.

Once a risk functional $\rho$ is fixed (Sections~\ref{sec:naive} to~\ref{sec:oce} below construct it), the quantity
\[
    \pi(0):=\inf_{\delta\in\HH}\rho\bigl(W(\delta)\bigr)
\]
measures the risk of the position of a trader with no liability who is free to trade. A proper definition and the systematic study of this optimization problem is the subject of Chapter~3. Following B\"uhler et al.\ \cite[Section~3.1]{buehler}, the market is called \emph{irrelevant} for the trader if $\pi(0)=-\infty$: by pure trading, the risk can be driven below every bound, and the hedging and pricing problem degenerates (every liability would have indifference price $-\infty$). Irrelevance is a joint property of the market \emph{and} the risk measure, and it can occur without classical arbitrage.

\begin{example}\label{ex:stat-arb}
    The following is a simplified version of the statistical-arbitrage
    illustration of B\"uhler et al.\ \cite[Section~3.1]{buehler}. We will show, that a market without arbitrage can still be irrelevant. Consider a
    single period, one asset, no costs and no constraints, and the risk
    measure $\rho(X)=-\E[X]$ of Section~\ref{sec:naive} below. Let the
    price increment be
    \[
        S_1-S_0=1/2+\varepsilon,
    \]
    where $\varepsilon$ takes the values $\pm1$ with probability $1/2$ each. There is no arbitrage: every nonzero position loses money with probability $1/2$. But the strategy $\delta_0=m$ has
    \[
        \rho\bigl(W(\delta)\bigr)=-\E[m(1/2+\varepsilon)]=-m/2\longrightarrow-\infty\qquad(m\to\infty),
    \]
    so $\pi(0)=-\infty$ leading to a irrelevent market. With a more risk averse $p$, the market could be relevant.
\end{example}

Relevance of the market is a condition on the position $X=0$ only. It is natural to ask whether it already guarantees feasibility, meaning $\pi(X)>-\infty$, for all other positions. On a finite probability space, B\"uhler et al.\ \cite[Corollary~3.3]{buehler} answer this affirmatively: since $\pi$ is monotone and cash-invariant (this is proved in Chapter~3), one has for every $X$
\[
    \pi(X)\ge\pi\bigl(\sup X\bigr)=\pi(0)-\sup X>-\infty,
\]
because on a finite $\Omega$ every random variable satisfies $\sup X<\infty$. On a general probability space this argument does not survive for unbounded $X$. This is the first appearance of the finite-versus-general theme of this thesis, and an early hint of why bounded liabilities, hypothesis (A1) of Chapter~3, play a distinguished role.

\section{Measuring risk: two na\"ive attempts}\label{sec:naive}

We now turn to the second half of the question: given the random P\&L \eqref{eq:PL}, how should one quantify its riskiness? We look for a functional $\rho$ that maps a random position $X$ to a number $\rho(X)$, interpreted as \emph{the amount of cash that must be added to the position to make it acceptable}. A positive $\rho(X)$ is a capital requirement; a negative $\rho(X)$ means cash can be withdrawn without making the position unacceptable. Before giving the axioms, we examine two natural first candidates and discuss why they fail to be good risk measures. Their failures are instructive because each one points directly at one of the axioms below.

\begin{example}\label{ex:neg-exp}
    \textbf{Negative expectation.}
    The map $\rho(X)=-\E[X]$ seems very natural in the cash interpretation: The exact amount of cash, that must be added to $X$ so that the expected P\&L becomes nonnegative, is also the amount that makes the position acceptable. Its defect is that it ignores the distribution of $X$ entirely: a sure gain of $1$ and a coin flip paying $+1000$ or $-998$ with probability $1/2$ each have the same expectation, hence the same risk number, though no risk-averse trader is indifferent between them. Example~\ref{ex:stat-arb} showed that under the negative expectation the market gets irrelevant at the slightest presence of statistical arbitrage. Whatever $\rho$ we choose, it must be able to distinguish positions by more than their mean.
\end{example}

\begin{example}\label{ex:variance}
    \textbf{Variance.}
    The map $\rho(X)=\mathrm{Var}(X)$ directly tackels the missing ingredient of the negative expectation: It punishes uncertainty. Furthermore, it is the canonical measure of risk in Markowitz portfolio theory \cite{markowitz1952}. Unfortunately, it still fails the cash interpretation twice. First, it penalizes gains and losses symmetrically: a position paying $+1000$ or $+998$ with probability $1/2$ has the same variance as one paying $-1000$ or $-998$, yet only the latter is risky in any economic sense. In particular the variance is not monotone: improving a position in every scenario can increase its risk number. Second, $\mathrm{Var}(X+c)=\mathrm{Var}(X)$ for every cash amount $c$, so adding cash does not reduce the risk number at all; the number cannot be a capital requirement, since no amount of capital ever makes the position acceptable. Both failures make economic nonsense of the reading of $\rho(X)$ as required capital.
\end{example}

These two failures motivate the axioms of the next section: monotonicity (better positions need less capital) and cash-invariance (added cash reduces the requirement one for one), together with convexity (diversification does not increase risk).

\section{Convex risk measures}\label{sec:crm}

A convex risk measure is a map on a vector space $\mathcal{X}$ of real-valued random variables on $(\Omega,\F,\PP)$. Before stating the axioms, we fix the domain. On a finite probability space with $N$ atoms, every real-valued random variable may be identified with a vector in $\R^N$, so the choice of $\mathcal{X}$ is redundant: one may take $\mathcal{X}=\R^N$, and any real-valued convex map $\rho\colon\R^N\to\R$ is automatically continuous. On a general probability space the choice matters. Taking $\mathcal{X}=L^\infty$ keeps $\rho$ real-valued under mild conditions, but excludes many positions that appear later (unbounded gains, unbounded liabilities). Following the needs of Chapter~3, we take the largest natural domain: $\mathcal{X}=L^0$, the space of all real-valued measurable random variables. On this space a useful risk measure may take the values $\pm\infty$, so we work with extended-real maps $\rho\colon L^0\to[-\infty,\infty]$. The axioms below are those of F\"ollmer and Schied \cite[Definition~4.1]{foellmerschied}, extended from a real-valued setting on a subspace of $L^\infty$ to this $L^0$-valued range; the same axioms are used by B\"uhler et al.\ \cite[Definition~1]{buehler}.

\begin{definition}\label{def:crm}
    A map $\rho\colon L^0\to [-\infty,\infty]$ is a \emph{convex risk measure} if for all $X,Y\in L^0$:
    \begin{enumerate}[leftmargin=*,label=\textup{(\roman*)}]
        \item \emph{monotonicity:} $X\le Y$ $\PP$-a.s.\ implies $\rho(X)\ge\rho(Y)$;
        \item \emph{cash-invariance:} $\rho(X+c)=\rho(X)-c$ for every $c\in\R$;
        \item \emph{convexity:} $\rho(\alpha X+(1-\alpha)Y)\le\alpha\rho(X)+(1-\alpha)\rho(Y)$ for every $\alpha\in[0,1]$.
    \end{enumerate}
\end{definition}

Each axiom has a direct economic reading, with the usual conventions of extended-real arithmetic when $\rho$ takes the values $\pm\infty$. Monotonicity says that a position that pays more in every scenario needs no more capital. Cash-invariance says that capital already in the position counts toward the requirement one for one. Also, it makes the unit of $\rho(X)$ the exact cash amount to render $X$ acceptable whenever $\rho(X)\in\R$
\[
    \rho\bigl(X+\rho(X)\bigr)=\rho(X)-\rho(X)=0.
\]
If $\rho(X)=+\infty$, no finite cash amount can push the position to be acceptable. Conversely, if $\rho(X)=-\infty$, the position is infinitely attractive and any finite cash amount can be withdrawn while keeping risk nonpositive. Convexity encodes the diversification principle: the risk of a mixture is at most the mixture of the risks, so splitting capital between two positions is never penalized. The negative expectation satisfies all three axioms (its failure in Section~\ref{sec:naive} is one of risk sensitivity, not of the axioms; it will be excluded later by the continuity hypothesis (A3) of Chapter~3), while the variance violates (i) and (ii). On $L^0$, convexity alone yields no automatic continuity: continuity of $\rho$ along the limits appearing in the proofs must be assumed (hypothesis (A3) of Chapter~3) or derived from structure (Lemma~4.4 of Chapter~4). % FAKE REF: lem:usc

\section{Optimized certainty equivalents}\label{sec:oce}

The convex risk measures used in this thesis all arise from a single construction, the \emph{optimized certainty equivalent} (OCE) of Ben-Tal and Teboulle \cite{bentalteboulle}, which B\"uhler et al.\ \cite[Section~3.3]{buehler} also take as their working class. The construction turns a one-dimensional loss function into a risk measure.

Throughout this section, $\ell\colon\R\to\R$ is a \emph{loss function}: convex, non-decreasing and finite on all of $\R$ (hence continuous). The value $\ell(x)$ is read as the disutility of a loss of size $x$; monotonicity says larger losses hurt more, convexity says they hurt disproportionately more.

\begin{definition}\label{def:oce}
    The \emph{optimized certainty equivalent} associated with $\ell$ is
    \begin{equation}\label{eq:OCE}
        \rho_\ell(X):=\inf_{w\in\R}\Bigl(w+\E\bigl[\ell(-X-w)\bigr]\Bigr),
    \end{equation}
    for all $X$ for which the expectation is well defined in $(-\infty,\infty]$ for every $w\in\R$.
\end{definition}

The interpretation of \eqref{eq:OCE} is a capital allocation: the trader sets aside the cash reserve $w$ today and suffers the disutility $\ell(-X-w)$ of the shortfall of the reserved position; optimizing over $w$ gives the best such allocation.

For \eqref{eq:OCE} to define a sensible risk measure, the loss function must be normalized. We adopt the normalization assumptions of Ben-Tal and Teboulle \cite[Definition~2.1, Equation~(2.1)]{bentalteboulle} by translating their utility function into our loss function via $u(x)=-\ell(-x)$.

\begin{definition}\label{def:subdifferential}
    The \emph{subdifferential} of the convex function $\ell$ at $0$ is the set
    \[
        \partial\ell(0):=\bigl\{s\in\R:\ell(x)\ge\ell(0)+sx\text{ for all }x\in\R\bigr\}.
    \]
    Elements of $\partial\ell(0)$ are called \emph{subgradients} of $\ell$ at $0$.
\end{definition}

\begin{assumption}\label{ass:normalization}
    The loss function $\ell$ satisfies $\ell(0)=0$ and $1\in\partial\ell(0)$.
\end{assumption}
If $\ell$ is differentiable at $0$, then $\partial\ell(0)=\{\ell'(0)\}$, so the condition $1\in\partial\ell(0)$ is equivalent to $\ell'(0)=1$.

The normalization $\ell(0)=0$ is a mere calibration of the zero level of disutility. Combined with $1\in\partial\ell(0)$
\[
    \ell(x)\ge \ell(0) + 1\cdot x=x\qquad\text{for all }x\in\R.
\]
Beyond this, the slope condition is essential for the following reason. If every subgradient of $\ell$ were strictly below $1$, say $\ell'\le\gamma<1$ everywhere, then for every fixed $x$,
\[
    w+\ell(-x-w)\le w+\ell(-x)+\gamma|w|\longrightarrow-\infty\qquad(w\to-\infty),
\]
so the infimum in \eqref{eq:OCE} would equal $-\infty$ for \emph{every} position $X$: the OCE would degenerate to $\rho_\ell\equiv-\infty$. Some slope of size at least $1$ is therefore necessary; placing it at $0$, together with $\ell(0)=0$, is the convenient normalization.

\begin{remark}\label{rem:oce-welldef}
    On a finite probability space, the expectation in \eqref{eq:OCE} is a finite sum and the definition needs no discussion. On a general $\Omega$, well-definedness must be checked. Let $X\in L^1$. Decompose $\ell(-X-w)$ into positive and negative parts. The negative part is controlled. By Assumption~\ref{ass:normalization} we have $\ell(x)\ge x$ which yields $\ell(-X-w)\ge-X-w$ and hence $\ell(-X-w)^-\le|X|+|w|\in L^1$. The expectation $\E[\ell(-X-w)]$ is therefore well defined in $(-\infty,\infty]$ for every $w$, and \eqref{eq:OCE} makes sense, with $\rho_\ell(X)=+\infty$ possible when the positive part fails to be integrable for every $w$. Moreover, the same inequality yields
    \[
        w+\E\bigl[\ell(-X-w)\bigr]\ge w+\E[-X-w]=-\E[X],
    \]
    so $\rho_\ell(X)\ge-\E[X]>-\infty$ for every $X\in L^1$: the OCE never takes the value $-\infty$ on $L^1$. On positions outside $L^1$ we use \eqref{eq:OCE} with the conventions of extended-real arithmetic, consistently with the domain choice in Section~\ref{sec:crm}.
\end{remark}

The next lemma is the general-$\Omega$ version of B\"uhler et al.\ \cite[Lemma~3.5]{buehler}; the proof of the three axioms follows theirs, rewritten on $L^0$ with extended-real expectations. Throughout, for a random variable $Y$ we write
\[
    \E[Y]:=\E[Y^+]-\E[Y^-]\in[-\infty,\infty],
\]
with the convention $(+\infty)-(+\infty):=+\infty$, so that \eqref{eq:OCE} defines a map $\rho_\ell\colon L^0\to[-\infty,\infty]$.

\begin{lemma}\label{lem:oce-crm}
    Under Assumption~\ref{ass:normalization}, the map $\rho_\ell$ defined by \eqref{eq:OCE} is a convex risk measure on $L^0$.
\end{lemma}

\begin{proof}
    We verify the three axioms of Definition~\ref{def:crm}, allowing the values $\pm\infty$ with the usual arithmetic $(\pm\infty)-c=\pm\infty$ for $c\in\R$ and $\alpha(\pm\infty)=(\pm\infty)$ for $\alpha\in(0,1]$.

    \emph{Monotonicity.} If $X_1\le X_2$ $\PP$-a.s., then $-X_1-w\ge-X_2-w$ $\PP$-a.s.\ for every $w$, so monotonicity of $\ell$ gives $\ell(-X_1-w)\ge\ell(-X_2-w)$ a.s. Taking extended-real expectations preserves the inequality, and taking the infimum over $w$ yields $\rho_\ell(X_1)\ge\rho_\ell(X_2)$.

    \emph{Cash-invariance.} For $c\in\R$, the substitution $w'=w+c$ gives
    \begin{align*}
        \rho_\ell(X+c) & =\inf_{w\in\R}\Bigl(w+\E\bigl[\ell(-X-c-w)\bigr]\Bigr)                    \\
                       & =\inf_{w'\in\R}\Bigl(w'-c+\E\bigl[\ell(-X-w')\bigr]\Bigr)=\rho_\ell(X)-c,
    \end{align*}
    including when $\rho_\ell(X)=\pm\infty$.

    \emph{Convexity.} Fix $\alpha\in[0,1]$, positions $X_1,X_2\in L^0$ and reserves $w_1,w_2\in\R$, and set $w_\alpha:=\alpha w_1+(1-\alpha)w_2$. Convexity of $\ell$ applied $\omega$-wise gives
    \[
        \ell\bigl(-\alpha X_1-(1-\alpha)X_2-w_\alpha\bigr)\le\alpha\ell(-X_1-w_1)+(1-\alpha)\ell(-X_2-w_2)\as
    \]
    Adding $w_\alpha$ and taking extended-real expectations,
    \begin{align*}
        w_\alpha+\E\bigl[\ell(-\alpha X_1-(1-\alpha)X_2-w_\alpha)\bigr]
         & \le\alpha\Bigl(w_1+\E\bigl[\ell(-X_1-w_1)\bigr]\Bigr)          \\
         & \qquad+(1-\alpha)\Bigl(w_2+\E\bigl[\ell(-X_2-w_2)\bigr]\Bigr),
    \end{align*}
    where the right-hand side is read in $[-\infty,\infty]$ (if it equals $+\infty$, the inequality is automatic; if it equals $-\infty$, both summands are $-\infty$ for $\alpha\in(0,1)$). The left-hand side is at least $\rho_\ell(\alpha X_1+(1-\alpha)X_2)$, so
    \begin{align*}
        \rho_\ell\bigl(\alpha X_1+(1-\alpha)X_2\bigr)
         & \le\alpha\Bigl(w_1+\E\bigl[\ell(-X_1-w_1)\bigr]\Bigr)          \\
         & \qquad+(1-\alpha)\Bigl(w_2+\E\bigl[\ell(-X_2-w_2)\bigr]\Bigr).
    \end{align*}
    Taking the infimum over $w_1$ and $w_2$ separately yields convexity.
\end{proof}

\begin{remark}\label{rem:oce-on-L1}
    Combining Remark~\ref{rem:oce-welldef} with Lemma~\ref{lem:oce-crm} recovers the $L^1$-statement that corresponds to B\"uhler et al.\ \cite[Lemma~3.5]{buehler}: the restriction of $\rho_\ell$ to $L^1$ is a convex risk measure with values in $(-\infty,\infty]$.
\end{remark}

We close the chapter with the two families of OCEs that serve as running examples in the rest of the thesis. They sit at opposite ends of a risk-aversion dial and demonstrate the flexibility of the class.

\begin{example}\label{ex:entropic}
    (Entropic risk measure.) For risk aversion $\lambda>0$, take the loss function
    \[
        \ell_\lambda(x):=\frac{e^{\lambda x}-1}{\lambda},
    \]
    which is convex, increasing and satisfies $\ell_\lambda(0)=0$ and $\ell_\lambda'(0)=1$, so Assumption~\ref{ass:normalization} holds. Fix $X$ with $\E[e^{-\lambda X}]<\infty$ and set
    \[
        f(w):=w+\E\bigl[\ell_\lambda(-X-w)\bigr]=w+\frac{1}{\lambda}\Bigl(e^{-\lambda w}\E\bigl[e^{-\lambda X}\bigr]-1\Bigr).
    \]
    The function $f$ is smooth and strictly convex with $f'(w)=1-e^{-\lambda w}\E[e^{-\lambda X}]$, so the unique minimizer is
    \[
        w^*=\frac{1}{\lambda}\log\E\bigl[e^{-\lambda X}\bigr],
    \]
    and since $e^{-\lambda w^*}\E[e^{-\lambda X}]=1$, one gets $f(w^*)=w^*$. Hence
    \[
        \rho_{\ell_\lambda}(X)=\frac{1}{\lambda}\log\E\bigl[e^{-\lambda X}\bigr],
    \]
    the \emph{entropic risk measure} with parameter $\lambda$; if $\E[e^{-\lambda X}]=\infty$, then $f\equiv\infty$ and $\rho_{\ell_\lambda}(X)=\infty$, consistently with the extended-real convention.
\end{example}

\begin{example}\label{ex:cvar}
    (Conditional value at risk / expected shortfall.) For a level $\alpha\in[0,1)$, take the loss function
    \[
        \ell_\alpha(x):=\frac{\max(x,0)}{1-\alpha},
    \]
    which is convex and non-decreasing with $\ell_\alpha(0)=0$ and $\partial\ell_\alpha(0)=[0,1/(1-\alpha)]\ni1$, so Assumption~\ref{ass:normalization} holds. The resulting OCE is the \emph{average value at risk} (also called conditional value at risk or expected shortfall) at level $1-\alpha$
    \[
        \rho_{\ell_\alpha}(X)=\inf_{w\in\R}\Bigl(w+\frac{1}{1-\alpha}\E\bigl[(-X-w)^+\bigr]\Bigr)=\mathrm{CVaR}_{1-\alpha}(X).
    \]
    It is the average of the losses that happen in the extreme tail with probability smaller than $1-\alpha$. Indeed: let $X\in L^1$ continuously distributed and write $L:=-X$ for the loss. The map $g(w):=w+\frac{1}{1-\alpha}\E[(L-w)^+]$ is convex in $w$ and, since $\PP[L=w]=0$ for every $w$, differentiable with
    \[
        g'(w)=1-\frac{1}{1-\alpha}\PP[L>w].
    \]
    The map $w\mapsto\PP[L>w]$ decreases continuously from $1$ to $0$, so the first-order condition $g'(w^*)=0$ has a solution $w^*$ with
    \[
        \PP[L>w^*]=1-\alpha,
    \]
    that is, $w^*$ is an $\alpha$-quantile of $L$, the value at risk of the position. Then
    \begin{align*}
        \rho_{\ell_\alpha}(X)=g(w^*) & =w^*+\frac{1}{1-\alpha}\E\bigl[(L-w^*)\mathbf{1}_{\{L>w^*\}}\bigr] \\
                                     & =w^*+\E\bigl[L-w^*\bigm|L>w^*\bigr]=\E\bigl[L\bigm|L>w^*\bigr]:
    \end{align*}
    the OCE equals the expected loss conditional on the tail event $\{L>w^*\}$ beyond the value at risk, an event of probability $1-\alpha$.
\end{example}

\begin{thebibliography}{9}

    \bibitem{bentalteboulle}
    A.~Ben-Tal and M.~Teboulle.
    An old-new concept of convex risk measures: the optimized certainty equivalent.
    \emph{Mathematical Finance} \textbf{17}(3), 449--476, 2007.

    \bibitem{buehler}
    H.~B\"uhler, L.~Gonon, J.~Teichmann, and B.~Wood.
    Deep hedging.
    \emph{Quantitative Finance} \textbf{19}(8), 1271--1291, 2019.

    \bibitem{foellmerschied}
    H.~F\"ollmer and A.~Schied.
    \emph{Stochastic Finance: An Introduction in Discrete Time}.
    4th edition, De Gruyter, Berlin, 2016.

    \bibitem{markowitz1952}
    H.~Markowitz.
    Portfolio selection.
    \emph{The Journal of Finance} \textbf{7}(1), 77--91, 1952.

\end{thebibliography}

\end{document}
